\documentclass[11pt,a4paper]{coverletter}

\senderblock
  {Radheem Bin Razi}
  {Distributed Systems \& Cloud-Native Engineer}
  {Ilmenau, Germany \textbar\ \href{mailto:sheikh.radheem@gmail.com}{sheikh.radheem@gmail.com} \textbar\ \href{https://radheem.github.io/my-cv}{portfolio}}

\begin{document}

%% ===== ENGLISH =====
\recipient{Aroundhome}{Hiring Team}{}

\opening{Dear Hiring Team,}

Helping new partners get from a first registration to an active partnership is the kind of problem where good backend engineering is felt directly by real people, so the Senior Software Engineer role on your partner-journey team was an easy one to get excited about. I am a backend and distributed-systems engineer with five years of experience in Go and Python, and I would bring hands-on system design, microservice delivery, and a strong reliability mindset to Aroundhome.

Most of my recent work has been designing event-driven microservices and owning them end to end. On a stealth distributed-systems platform I built Go services with gRPC and Protobuf-first APIs, and led the migration off Dapr to native NATS JetStream and KV — exactly the kind of architecture modernization you describe, replacing a heavier abstraction with something simpler and more observable. Earlier, as a design authority at Bluefin Exchange, I reviewed code and designs for the core services of a high-throughput crypto exchange and drove improvements where the system was weakest. I care about what happens after release, too: my projects wire up OpenTelemetry, Prometheus, and Grafana so reliability work and alerting are grounded in real signals rather than guesses.

I also work the way your team does. I default to clean, tested code and treat code reviews and pair programming as how a team gets better, not as overhead — I have led and mentored junior engineers throughout my career. And I lean on AI heavily in my own development: my Second Brain project is a self-hosted RAG service with local LLM inference and pgvector search, and I have built MCP tooling that exposes a platform to LLM agents. Turning that into faster understanding of a codebase and quicker delivery is second nature to me.

I am available full time and based in Germany, so a Berlin-hybrid or remote arrangement fits well, and I can relocate within Germany within a few weeks. I am already authorized to work here and am Blue Card–ready — for a qualifying offer you would only need to issue the contract, with no sponsorship overhead. Thank you for your time and consideration; I would welcome the chance to talk about how I can help your partner-journey team.

\closing{Sincerely,}{Radheem Bin Razi}

\clearpage
\selectlanguage{ngerman}

%% ===== DEUTSCH =====
\recipient{Aroundhome}{Personalabteilung}{}

\opening{Sehr geehrtes Hiring-Team,}

Neuen Partnern dabei zu helfen, von der ersten Registrierung zu einer aktiven Partnerschaft zu gelangen, ist genau die Art von Problem, bei der gute Backend-Entwicklung von echten Menschen unmittelbar spürbar wird — die Rolle als Senior Software Engineer in Ihrem Partner-Journey-Team hat mich daher sofort begeistert. Ich bin Backend- und Distributed-Systems-Engineer mit fünf Jahren Erfahrung in Go und Python und würde praxiserprobtes Systemdesign, die Auslieferung von Microservices und ein ausgeprägtes Zuverlässigkeitsdenken zu Aroundhome mitbringen.

Ein Großteil meiner jüngsten Arbeit bestand darin, ereignisgesteuerte Microservices zu entwerfen und sie end-to-end zu verantworten. Auf einer Distributed-Systems-Plattform im Stealth-Modus habe ich Go-Dienste mit gRPC und Protobuf-First-APIs entwickelt und die Migration von Dapr hin zu nativem NATS JetStream und KV geleitet — genau die Art von Architekturmodernisierung, die Sie beschreiben, bei der eine schwergewichtige Abstraktion durch etwas Einfacheres und besser Beobachtbares ersetzt wird. Zuvor habe ich als Design-Authority bei Bluefin Exchange Code und Designs für die Kerndienste einer durchsatzstarken Krypto-Börse geprüft und Verbesserungen dort vorangetrieben, wo das System am schwächsten war. Mir ist auch wichtig, was nach dem Release passiert: Meine Projekte binden OpenTelemetry, Prometheus und Grafana ein, sodass Zuverlässigkeitsarbeit und Alerting auf realen Signalen statt auf Vermutungen beruhen.

Ich arbeite zudem so, wie es Ihr Team tut. Ich setze standardmäßig auf sauberen, getesteten Code und betrachte Code-Reviews und Pair Programming als die Art, wie ein Team besser wird, nicht als Mehraufwand — ich habe im Laufe meiner Karriere durchgehend Junior-Engineers geleitet und betreut. Und ich nutze KI in meiner eigenen Entwicklung intensiv: Mein Second-Brain-Projekt ist ein selbst gehosteter RAG-Dienst mit lokaler LLM-Inferenz und pgvector-Suche, und ich habe MCP-Tooling entwickelt, das eine Plattform für LLM-Agenten zugänglich macht. Daraus ein schnelleres Verständnis einer Codebasis und eine zügigere Auslieferung zu machen, ist für mich selbstverständlich.

Ich bin in Vollzeit verfügbar und in Deutschland ansässig, sodass ein Berlin-Hybrid- oder Remote-Modell gut passt, und ich kann innerhalb weniger Wochen innerhalb Deutschlands umziehen. Ich bin bereits arbeitsberechtigt und Blue-Card-bereit — für ein qualifizierendes Angebot müssten Sie lediglich den Vertrag ausstellen, ohne jeglichen Aufwand für ein Sponsoring. Vielen Dank für Ihre Zeit und Ihr Interesse; ich würde mich über die Gelegenheit freuen, darüber zu sprechen, wie ich Ihr Partner-Journey-Team unterstützen kann.

\closing{Mit freundlichen Grüßen,}{Radheem Bin Razi}

\end{document}
